% !TEX root = ./LP-main.tex

We start this chapter by introducing a standard textbook-style linear program,~(\ref{P0}). We then investigate the celebrated Farkas' Lemma, Theorem~\ref{farkas}, to give a first characterization for the feasibility of a linear program. Farkas' Lemma is sufficient for proving the strong duality of a linear program; therefore, we prove a specific case for the minimax theorem in Theorem~\ref{minimax}; we then generalize it to a cone-version of the statement, in which the proof is straightforward. This then provides enough motivation to properly define the dual of a primal problem. 

In the second part of the chapter, we move to a geometric point of view for Farkas' Lemma, Theorem~\ref{farkas geo}. The two versions of Farkas' Lemma are also known as the H-representation and V-representation of a polyhedral set, which we will then further investigate in Chapter~\ref{chapter: polyhedral geometry}. 

\section{Farkas' Lemma and Dual Problem}

In linear optimization problems, the constraints are expressed as linear transformations; we may exploit the bounds in inner product spaces to obtain an upper bound on the maximum. We first restate the optimization problem. Consider the system of linear inequalities and the optimization problem that
\begin{equation}\label{P0}\tag{P0}
\begin{array}{ll}
\mathrm{min}_{\bmx\in\rr^n}  \quad & \inner{\bmc, \bmx}\\
\text{subject to} \quad & M\bmx  \leq \bmb.
\end{array}
\end{equation}
and we may view $M$ as a linear transformation $M:\rr^n \to \rr^m$; in a way, the optimization problem (P0) is to study the row space of the matrix $M$, where one needs to find a vector in $\rr^n$ so that it has a direction close to the objective $\bmc$ and obeys the constraint $M$.  

Before looking for solutions to the linear program (\ref{P0}), we first want to check whether solutions exist. The following theorem shows that the feasibility of a linear system is completely determined by an associated dual system in $\rr^m$. Recall that the hyperplane separation theorem (Theorem~\ref{prop: strict hyperplane separation}) states the following: if a point $\bmb$ does not lie inside a closed convex set $Z$, you can always draw a hyperplane that completely separates $\bmb$ from $Z$ and does not touch $\bmb$ at all. 


\begin{theorem}[Farkas' Lemma]\label{farkas}
Let $M\in\rr^{m\times n}, \bmb\in\rr^m$ and $x_1,$ $x_2$, $\dots$, $x_n$, $u_1$, $u_2$, $\dots$, $u_m$ be a list of variables written in their vector forms $\bmx\in\rr^n$ and $\bmu\in\rr^m$. Consider the following two systems of linear inequalities. 
\begin{equation}\label{eqn: P}
M\bmx\leq b
\end{equation} and 
\begin{equation}\label{eqn: D}
\begin{array}{rl}
M^\top \bmu & = 0\\
\bmu & \geq 0\\
\inner{\bmu, \bmb} & < 0. \\
\end{array}
\end{equation}
The system of inequalities (\ref{eqn: P}) has a feasible solution if and only if the system of equalities and inequalities (\ref{eqn: D}) has no solution. 
\end{theorem}
\begin{proof}
The forward direction, suppose that there are $\bmx$ and $\bmu$ so that (\ref{eqn: P}) and (\ref{eqn: D}) both have solutions; we have that
\[ 0 > \inner{\bmu, \bmb} \geq \inner{\bmu, M\bmx} = \inner{M^\top \bmu , \bmx} = 0, \] which leads to a contradiction and, thus, (\ref{eqn: D}) has no solution. For the converse direction, we proceed via contrapositive, supposing (\ref{eqn: P}) has no solution; we will show that (\ref{eqn: D}) has a solution. 

The assumption that (\ref{eqn: P}) has no solution means the set $\{M\bmx: \bmx\in\rr^n\}$ does not include $\bmb$; since, otherwise, we would have a solution $\bmx\rast$ so that $M\bmx\rast = \bmb \leq \bmb$. In particular, we have that $Z = \{M\bmx + \bm s:  \bmx\in\rr^n, \bms\in\rr_+^m\}$ does not include $\bmb$. By the {\em hyperplane separation theorem}, there exists $\bmu^\ast\in\rr^m$ and $q\in\rr$ such that $\inner{\bmu\rast, \bmb} < q$ and $\inner{\bmu\rast, \bmz} \geq q$ for every $\bmz\in Z$. In particular, since $M$ is a linear transformation, $\bmx$ is taken over $\rr^n$, $\bmb\not\in \operatorname{Im}(M) := \{\bmy\in\rr^m: \bmy= M\bmx\}$, and $\operatorname{Im}(M)$ is a subspace of $\rr^m$, we have that $Z$ is a convex cone and we may take $q=0$. 

We now verify that $\bmu\rast$ is a solution to (\ref{eqn: D}). First, we have $\inner{\bmu\rast, \bmb} < 0$ by the choice according to the separation theorem. By the construction of $Z$, we may say for all $\bmx\in\rr^n$ and $\bms\in\rr_+^m$, 
\begin{equation}\label{eqn: construction of Z}\inner{\bmu\rast, M\bmx + \bm s} = \inner{\bmu\rast, M\bmx} + \inner{\bmu\rast, \bm s} \geq 0,
\end{equation} 
where the inequality is from the separation theorem. When we fix $x = 0$, Inequality (\ref{eqn: construction of Z}) gives that $\inner{\bmu\rast, \bm s}\geq 0$ for any $\bm s\in\rr_+^m$: we may conclude that $u\rast\in\rr_+^m$ (i.e., $\bmy\rast\geq 0$). 

Finally, we assume $\bms = 0$, we have that $ \inner{\bm u\rast, M\bmx} + \inner{\bmu\rast, \bm s} =  \inner{\bmu\rast, M\bmx} = \inner{M^\top \bmu\rast, \bmx} \geq 0$. Now, by properties of inner products, for every $\bmx\in\rr^n$, we have that $\inner{M^\top \bmu\rast, -x} = -\inner{M^\top \bmu\rast, \bmx}\leq 0$. This then implies that $\inner{M^\top \bmu\rast, \bmx} = 0$, as it is the only number that equals its own negation. By the positive-semidefiniteness of inner products, it follows that $M^\top \bmu\rast = 0$. 
\end{proof}

Observe that, we want only the solution where $M\bmx\leq \bmb$ in (\ref{P0}), which is equivalent to $M\bmx - \bmb\in\rr^n_-$ and, otherwise, discard may as $-\infty$. Consider the following notion to reformulate the primal problem (\ref{P0}) in a more analytical form.  
The {\em geometrical indicator function} is denoted and defined as $\delta_S: \rr^m \to \rr\cup\{\infty\}$ where
$$ \delta_S(\bms) = \left\{ \begin{array}{ll} 
 0 & \text{ if } \bms\in S, \text{ and }\\
\infty & \text{ otherwise.} 
\end{array} \right. $$
The subscript is sometimes replaced by an event with respect to the input; for example, if $S = \{\bmx\in \rr^n: M \bmx\leq \bmb\}$ for some $M \in\rr^{m\times n}$ and $\bmb\in \rr^m$, then we may write $\delta_S(\bmx_0) = \delta_{M\bmx_0\leq \bmb}(\bmx_0).$ We now reformulate the program (\ref{P0}) as 
\begin{equation}\label{a}
 z\rast = \max\{\inner{\bmc, \bmx}:  M \bmx\leq \bmb, \xx\in \rr^n\} = \sup_{\bmx\in\rr^n}\left\{ \inner{\bmc, \xx} - \delta_{\rr^m_-}( M \bm x - \bmb)\right\} 
 \end{equation}  
where $z\rast\in\rr\cup\{-\infty\}$.  Recall that, given a cone $S\subseteq V$, the polar cone is denoted and defined as $S^\circ = \{\bmy\in V: \forall \bmx\in S, \inner{\bmx, \bmy}\leq 0\}$. 



\begin{proposition}[Support function Representation via Fenchel-Moreau]\label{prop: dual representation}
Suppose that $S$ is a closed convex cone. We have that 
$\delta_{S}(x) = \sup\{\inner{x, u}:{u\in S^\circ}\}.$\hfill \qedsymbol
\end{proposition} 

Proposition~\ref{prop: dual representation} then rewrites the program defined in~(\ref{a}) as
\begin{align*}
 z\rast = & \sup_{\bmx\in\rr^n}\left\{ \inner{\bmc, \xx} -  \sup_{\bmu\in\rr^m_+}\left\{ \inner{\bmu, M\bm x - \bmb}  \right\}\right\} \\
 = & \sup_{\bmx\in\rr^n}\left\{ \inner{\bmc, \xx} +  \inf_{\bmu\in\rr^m_+}\left\{ \inner{\bmu, \bmb}-\inner{\bmu, M\bm x}  \right\}\right\}, \\ 
\end{align*}
which finally gets to 
\begin{equation}\label{eqn: z sup inf}
 z\rast =  \sup_{\bmx\in\rr^n} \inf_{\bmu\in\rr^m_+} \big\{ \inner{\bmc, \xx} + \inner{\bmu, \bmb}-\inner{\bmu, M\bm x}  \big\}. 
\end{equation}
Given $\bmb \in \rr^n$ and $\bmc \in \rr^m$, we may convienently define bivariate function $L$ over a subset of $\rr^n\times \rr^m$ where $L(\bmx, \bmu)  =  \inner{\bmc, \xx} +  \inner{\bmu, \bmb}-\inner{\bmu, M\bm x}$. In the aforementioned definition for $z\rast$, its domain is $\rr^n\times \rr^m_+$, which then gives $z\rast = \sup_{\bmx\in\rr^n} \inf_{\bmu\in\rr^m_+} L(\bm x, \bmu)$. Now, we consider the following sup-inf exchange. In particular, the following theorem is the fundamental theorem that establishes the connection between the solution of a linear optimization problem and that of its dual. 

\begin{theorem}\label{minimax}
Let $M\in \rr^{m\times n}, \bbb\in\rr^m, \bcc\in\rr^n$. Suppose that $L$ is a bivariate functional defined on $\rr^n\times \rr_+^{m}$ such that $L(\bmx, \bmu)  =  \inner{\bmc, \xx} +  \inner{\bmu, \bmb}-\inner{\bmu, M\bm x}$. If there is at least one feasible $\bmx\rast$ for $M\bmx\rast \leq \bmb$ and the supremum $z^* :=  \sup_{\bmx\in\rr^n} \inf_{\bmu\in\rr^m_+} L(\bm x, \bmu)$ is finite, then 
\begin{equation}\label{eqn: sup inf exchange}  \sup_{\bmx\in\rr^n} \inf_{\bmu\in\rr^m_+} L(\bm x, \bmu) =  \inf_{\bmu\in\rr^m_+} \sup_{\bmx\in\rr^n} L(\bm x, \bmu). \end{equation}
\end{theorem}
\begin{proof}
Let $X = \rr^n$ and $U =  \rr^m_+$. We first prove that \[ \sup_{\bmx\in X} \inf_{\bmu\in U} L(\bmx, \bmu) \leq \inf_{\bmu\in U} \sup_{\bmx\in X} L(\bmx, \bmu),\] purely by the definitions of supremum and infimum. For all $\bmx_0\in X$ and $\bmu_0\in U$, the following inequality holds by definition 
\[ \inf_{\bmu\in U} L(\bmx_0, \bmu) \leq L(\bmx_0, \bmu_0)\leq \sup_{\bmx\in X}L(\bmx, \bmu_0). \]
Following precisely from the definition of supremum, the condition that for all $\bmx_0\in X$, $ \inf_{\bmu\in U} L(\bmx_0, \bmu) \leq \sup_{\bmx\in X}L(\bmx, \bmu_0)$ provides the following inequality 
\[ \sup_{\bmx\in X}\inf_{\bmu\in U} L(\bmx, \bmu) \leq \sup_{\bmx\in X}L(\bmx, \bmu_0). \]
Again, since $\bmu_0\in U$ is chosen arbitrarily, we have that 
\[ \sup_{\bmx\in X}\inf_{\bmu\in U} L(\bmx, \bmu) \leq \inf_{\bmu\in U}\sup_{\bmx\in X}L(\bmx, \bmu), \]
which proves one direction. 

We now turn to the other direction that $ \sup_{\bmx\in X} \inf_{\bmu\in U} L(\bmx, \bmu) \geq \inf_{\bmu\in U} \sup_{\bmx\in X} L(\bmx, \bmu)$; toward a contradiction, suppose that $ \sup_{\bmx\in X} \inf_{\bmu\in U} L(\bmx, \bmu) < \inf_{\bmu\in U} \sup_{\bmx\in X} L(\bmx, \bmu)$. Observe that the polar dual of $\rr^n$ is $(\rr^n)^\circ = \{0\}$. We may apply the support function representation so that, for a fixed $\bmu_0$, we have  
\begin{align*}
\sup_{x\in X} L(\bmx, \bmu_0) = & \sup_{\bmx\in X}\left\{\inner{\bmc, \bmx} + \inner{\bmu_0, \bmb}-\inner{\bmu_0, M \bmx}\right\} \\
	= 	& \inner{\bmu_0, \bmb} + \sup_{\bmx\in \rr^n}\left\{ \inner{\bmc, \bmx}-\inner{\bmu, M \bmx}\right\} \\ 
	= 	& \inner{\bmu_0, \bmb} + \sup_{\bmx\in \rr^n}\left\{ \inner{\bmc - M^\top \bmu_0, \bmx}\right\} \\ 
	= 	& \inner{\bmu_0, \bmb} + \delta_{\{\bmu: \bmc- M^\top \bmu = 0\} }(\bmu_0). \\ 
 \end{align*}
By the contradiction hypothesis and the definition of an infimum, fix any $\bmu_0\in \rr_+^m$, we have the following chain of inequalities that
\[ \sup_{\bmx\in \rr^n} \inf_{\bmu\in \rr^m_+} L(\bmx, \bmu) < \inf_{\bmu\in \rr^m_+} \sup_{\bmx\in X} L(\bmx, \bmu) \leq \inner{\bmu_0, \bmb} + \delta_{\{\bmu: M^\top \bmu = c\}}(\bmu_0), \]
and $ \inner{\bmu_0, \bmb} + \delta_{\{\bmu: M^\top \bmu = \bmc\}}(\bmu_0) \in\{ \inner{\bmu_0, \bmb}, \infty\}. $ Thus, as we verified the two sides of the strict inequality and by the density of the real numbers, we are ensured that there is an $\alpha\in\rr$ so that 
\[ \max_{\bmx\in\rr^n}\{\inner{\bmc, \bmx}: M\bmx\leq \bbb\} =:  z\rast < \alpha < \inf_{\bmu\in \rr^m_+} \sup_{\bmx\in X} L(\bmx, \bmu).  \]

\noindent Since $z\rast$ is the supremum, that is, $\inner{\bmc, \bmx_0} < \alpha$ for all $\bmx_0\in\rr^n$, the following system does not have a solution
$$ \begin{bmatrix} M\\ -\bmc^\top \end{bmatrix} \bmx \leq \begin{bmatrix} \bm b\\ -\alpha\end{bmatrix}. $$
By Theorem~\ref{farkas}, let $\bmu = (u_1, u_2, \dots, u_m)^\top$ and $u_c$ be variables, the following systems of inequalities  
\begin{equation}\label{inhom} \begin{split}
\begin{bmatrix} M^\top\mid -\bmc \end{bmatrix}\begin{bmatrix} \bmu\\ u_c \end{bmatrix} = \bm0 \\
\begin{bmatrix} \bmu\\ u_c \end{bmatrix}\geq\bm0 \\
\inner{\begin{bmatrix} \bmu\\ u_c \end{bmatrix}, \begin{bmatrix} \bmb\\ -\alpha \end{bmatrix}} < 0 
\end{split}\end{equation}
has a solution. Let $\bmu\rast$ and $u_c\rast$ be a solution for the system (\ref{inhom}), and reformulate so we have 
\begin{equation}\label{eqn: aux D}  M^\top \bmu\rast =  u_c\rast\bmc, \quad \bmu\rast \geq\bm0,\quad u_c\rast \geq0,\quad \inner{\bmu\rast, \bmb} < \alpha  u_c\rast. \end{equation}

We first assume $u_c\rast\neq 0$, which (\ref{eqn: aux D}) gives that $u_c\rast > 0$. Let $\widetilde{\bmu} = \bmu\rast / u_c\rast$; since $u_c\rast\geq 0$, we obtain that 
$$  M^\top \widetilde\bmu =  \bmc, \quad 
\widetilde\bmu = \frac{\bmu\rast}{u_c\rast} \geq \bm0,\quad 
\inner{\widetilde\bmu, \bmb} = \frac{1}{u_c\rast}\inner{\bmu\rast, \bmb} < \alpha . $$  
However, we now have $\delta_{M^\top\widetilde\bmu = \bmc} (\widetilde\bmu) = 0$, which leads to the contradiction
\[ \inner{\widetilde\bmu, \bmb}  < \alpha< \inf_{u} \sup_{x} L( x, u) \leq \inner{\widetilde\bmu, \bmb}, \] where $\alpha$ is assumed to be strictly less than the infimum and the infimum is the lower bound for $\inner{\wt{\bmw}, \bmb}$ by definition.  

Alternatively, we assume $u_c\rast = 0$ and plug in to inequalities (\ref{eqn: aux D}) and
$$  M^\top \bmu\rast =  \bm0, \qquad \bmu\rast \geq \bm0,\qquad \inner{\bmu\rast, \bmb} < 0. $$
As we assumed its existence, let $\bmx_0$ be a feasible solution to the system of inequalities $M\bmx\leq\bmb$. Since $M^\top \bmu\rast =  \bm0$ and $\bmu\geq\bm0$, we now have a contradiction that 
$$ 0 = \inner{M^\top\bmu\rast, \bmx_0} = \inner{\bmu\rast, M \bmx_0} \leq \inner{\bmu\rast, \bmb} < 0. $$
This completes the proof. 
\end{proof}



We observe from the proof of Theorem~\ref{farkas} that, the change of the constraints and bounds in the primal problem (\ref{eqn: P}) leads to the change in the constraints in its dual problem (\ref{eqn: D}). 
\begin{itemize}
\item In (\ref{eqn: D}), the constraint $A^\top \bmy = \bm0$ is determined by the fact that $\bmx$ is taken from the entire space of $\rr^n$. If we instead restrict $\bmx\in\rr^n_+$ and change the proof accordingly, then the dual constraints in (\ref{eqn: D}) change to $A^\top \bmy \geq \bm0$. Indeed, in the proof, when the slack variable satisfies $\bms = \bm0$, then $\inner{\bmy\rast, A\bmx} = \inner{A^\top \bmy\rast, \bmx}\geq 0$ for all $\bmx\geq \bm0$ if and only if $A^\top \bmy\geq \bm0$. More generally, for any closed convex cone $C$, if $\bmx\in C$, then the dual condition becomes 
\[ M^\top \bmu\in C\rast\text{ where } C\rast := \{\bmy \in\rr^n: \forall \bmx\in C, \inner{\bmx, \bmy}\geq 0\}. \]
where $C\rast$ is known as the {\em dual cone} of $C$. 

\item The non-negativeness of $\bmu\in\rr^m_+$ is independent of the primal cone $C$ of where $\bmx$ is choosen. But dictated by the inequality direction $M \bmx\leq \bmb$. Observe that if we instead have $M\bmx\geq \bmb$, then we would take $Z = \{A\bmx : \bmx\in R^n\} + R^m_-$, which then gives $\bmy\leq\bm0$.  
%  determined by the construction of $Z = \{Ax: x\in X, x=0\} + \rr^m_+$, regardless of what $ X$ is. 
\item In a similar vein, the bivariate Lagrangian function in Theorem~\ref{minimax} can be defined more generally.
\end{itemize}
\begin{theorem}[Conic Minimax Theorem]\label{conic minimax}
Let $M\in \rr^{m\times n}, \bmb\in\rr^m, \bmc\in\rr^n$. Let $X\subset \rr^n$ and $U\subset\rr^m$ be closed convex cones and the bivariate Lagrangian $L: X\times U\to \rr$ is defined as $L(\bmx, \bmu)  =  \inner{\bmc, \xx} +  \inner{\bmu, \bmb}-\inner{\bmu, M\bm x}$. If there is at least one feasible $\bmx\rast\in X$ so that $\bmb - M\bmx\rast \in U $ and the supremum $z^* :=  \sup_{\bmx\in X} \inf_{\bmu\in U\rast} L(\bm x, \bmu)$ is finite, then 
\begin{equation} \tag*{\qedsymbol} \sup_{\bmx\in X} \inf_{\bmu\in U\rast} L(\bm x, \bmu) =  \inf_{\bmu\in U\rast} \sup_{\bmx\in X} L(\bm x, \bmu). \end{equation}
\end{theorem}  

Finally, note that the system~(\ref{inhom}) is precisely how to transform an optimization problem with an inhomogeneous system into a homogeneous system with no optimization tasks. As such, and along with the remarks above, we have the formal and complete form of a linear program as follows. Let $A\in\rr^{m\times n}$, $\bmb\in\rr^m$, and $\bmc\in\rr^n$; let $L\subset \rr^m$ and $K\subset \rr^n$ be two closed convex cones. Let $\bmx\in\rr^n$ and $\bmy\in\rr^m$ be two vectors of variables. When we define the Linear Problem (GP) as 
\begin{equation}\label{GP}\tag{GP}
\begin{array}{ll}
\mathrm{max}_{\bmx}  \quad & \inner{\bmc, \bmx}\\
\text{subject to} \quad & \bmb - A\bmx  \in L, \\
& \bmx\in K; 
\end{array}
\end{equation}
the {\em dual} problem of (\ref{GP}) follows and defines the Linear Problem (\ref{GD}) as
 \begin{equation}\label{GD}\tag{GD}
\begin{array}{ll}
\mathrm{min}_{\bmy}  \quad & \inner{\bmy, \bmb}\\
\text{subject to} \quad & A^\top\bmy  - \bmc  \in K\rast, \\
& \bmy\in L\rast.
\end{array} 
\end{equation}
If we choose $L = \rr^m_+$ and $K = \rr^n_+$, then we can define a primal problem (\ref{P}) such that 
\begin{equation}\label{P}\tag{P}
\begin{array}{ll}
\mathrm{max}_{\bmx\in\rr^n}  \quad & \inner{\bmc, \bmx}\\
\text{subject to} \quad & A\bmx  \leq \bmb, \\
& \bmx\geq 0, 
\end{array}
\end{equation}
and its dual problem (\ref{D}) as 
 \begin{equation}\label{D}\tag{D}
\begin{array}{ll}
\mathrm{min}_{\bmy\in\rr^m}  \quad & \inner{\bmy, \bmb}\\
\text{subject to} \quad & A^\top\bmy  \geq \bmc, \\
& \bmy\geq 0.
\end{array} 
\end{equation}
In particular, (P) has a solution if and only if (D) has a solution and, if (P) has a solution, then the maximum it attains corresponds to a solution to (D) that attains the same value as the minimum. However, there is a caveat: even if we are given an optimal solution to a primal problem, we do not immediately obtain an optimal solution for its dual problem, as the strong duality only guarantees the existence of an optimal solution, not an explicit construction. A relevant constructive solution will be discussed in Chapter~\ref{chapter: linear programming}. 

\section{Geometric Point of View for Dual Problems}

In Theorem~\ref{farkas}, we considered the existential argument for the feasibility of a system $M\bmx\leq \bmb$. We now turn to a constructive description of Farkas' Lemma. In particular, we state it in a less general form when we also consider the equality constraints; yet, an equality constraint $ ax = b$ can be expressed as two inequalities $ax\leq b$ and $-ax\geq b$ and, thus, equality constraints are not considered special cases for systems of linear inequalities.  

\begin{theorem}[Geometric Farkas' Lemma]\label{farkas geo}
Let $m,  d_E, d_A$ be positive integers and $E\in\rr^{m\times d_E}, A\in\rr^{m\times d_A}$, and $Y_{E, A} = \{\bmy\in \rr^{m}: E^\top \bm y = \bm0, A^\top \bm y \leq \bm0\}$. A vector $ \bmb\in\rr^{m}$ satisfies $\inner{\bm y, \bm b}\leq 0$ for every $\bmy\in Y_{E, A}$ if and only if there exists $\bm\lambda\in\rr^{d_E}$ and $\bm\mu\in\rr^{d_A}_+$ so that 
\begin{equation}\label{eqn: farkas geo} \bmb = E\bm\lambda + A\bm\mu. \end{equation}
\end{theorem} 
\begin{proof}
First, we reformulate the constraint (\ref{eqn: farkas geo}) into the following system for $\bm\lambda \in\rr^{d_E}$ and $\bm\mu \in\rr^{d_A}$ such that 
\begin{equation}\label{eqn: farkas geo aux}
\begin{bmatrix}
E & A \\
-E & - A \\
\bm O& \bm -I_{d_A}
\end{bmatrix}
\begin{bmatrix}
\bm\lambda \\ \bm\mu
\end{bmatrix}
\leq 
\begin{bmatrix}
\bmb\\ - \bmb\\ \bm 0
\end{bmatrix}, 
\end{equation}
where $\bm O$ denotes an all-zero matrix with appropriate dimensions. The two blocks $E\bm\lambda + A\bm\mu \leq \bmb$  and $- E\bm\lambda - A\bm\mu \leq - \bmb$ corresponds to the equality, and the block $\bm O\bm\lambda - I \bm\mu \leq \bm0$ corresponds to that $\bm\mu\geq\bm0$. Since we will apply Theorem~\ref{farkas}, consider the following dual system of equalities and inequalities with a list of variables ${\bmu}\in\rr^{m + m + d_A}$: 
\begin{equation}\begin{split}\label{eqn: geo farkas dual}
%constraint 1
\left[ \begin{array}{c|c|c}
E^\top  & -E^\top & \bm O\\
A^\top  & -A^\top  & -I_{d_A}
\end{array} \right]
{\bmu} & =  \bm0, \\
% constraint 2
{\bmu} \geq \bm0,  \\
% constraint 3
\inner{{\bmu}, 
\begin{bmatrix} 
\bmb\\  -\bmb\\  \bm0 
\end{bmatrix}
} & < 0\\
\end{split}\end{equation}

We first prove the forward implication; suppose $\inner{\bmy, \bmb} \leq 0$ for every $\bmy \in Y_{E, A}$ and, toward a contradiction, that the system~(\ref{eqn: farkas geo aux}) has no solution; by Theorem~\ref{farkas}, there is a solution $\tilde{\bmu}\in\rr^{m + m + d_A}$ for (\ref{eqn: geo farkas dual}). 
 We partition the vector 
\[\tilde{\bm u} = \begin{bmatrix}\bm y_1\\ \bm y_2\\ \bm y_3 \end{bmatrix}, \]
where $\bmy_1, \bmy_2\in\rr^{m}_+$ and $\bmy_3 \in\rr_+^{d_A}$ and the non-negativeness is given that that $\tilde\bmu\geq\bm0$, so that $E^\top \bmy_1 - E^\top \bmy_2 = \bm 0 $  and $A^\top\bm y_1 - A^\top\bm y_2 -  \bm y_3 = \bm 0$. Equivalently, we write that 
\[E^\top(\bmy_1- \bmy_2) = \bm0, \quad A^\top(\bmy_1 - \bmy_2) =  \bmy_3 =  \bm0, \quad \text{ and }\inner{\bmy_1, \bm b} - \inner{\bmy_2, \bmb} < 0, \]
where the last inequality is due to $\inner{\bmy_1, \bmb} + \inner{\bmy_2, -\bmb} + \inner{\bmy_3, \bm0} < 0$ as in (\ref{eqn: geo farkas dual}). Now, we set $\bmy_0 = \bmy_2 - \bmy_1$ and we have that $E^\top \bmy_0 = \bm 0$, $A^\top \bm y_0\leq \bm 0$; this means $\bmy_0\in Y_{E, A}$. Finally, $\inner{\bmy_0, \bmb} > 0$ contradicts the premise that $\inner{\bmy,\bmb}\leq 0$ for all $\bmy\in Y_{E, A}$. 

For the converse direction, suppose that $\bmb = E\bm\lambda_0 + A \bm\mu_0$ for some $\bm\lambda_0\in\rr^{d_E}$, $\bm\mu_0\in\rr^{d_A}_+$. We assume that $Y_{E, A}\neq\emptyset$ or, otherwise, we are done. By the existence of $\bm\lambda_0$ and $\bm\mu_0$, the system~(\ref{eqn: farkas geo aux}) has a solution and, by Theorem~\ref{farkas}, the system~(\ref{eqn: geo farkas dual}) has no solution. In the following, we fix a $\bmy\in Y_{E, A}$ and we will show that $\inner{\bmy, \bmb}\leq 0$.

For $i\in\{1,2,\dots, m\}$, let $y_i$ denote the $i$-th component of $\bmy$; we define $\bmu^+, \bmu^-\in\rr^m$ where each of their $i$-th component has that 
\[ u_i^+ = \left\{\begin{array}{ll} y_i & \text{ if } y_i\geq 0\\ 0 & \text{otherwise},
 \end{array} \right.
\qquad\text{ and }\qquad
 u_i^- = \left\{\begin{array}{ll} -y_i & \text{ if } y_i < 0 \\ 0 & \text{otherwise, }
  \end{array} \right.
\]
so that $\bmy = \bmu^+ - \bmu^-$. We now set 
\[{\bmu\rast} = \begin{bmatrix}\bmu^-\\ \bmu^+\\ - A^\top \bmy \end{bmatrix}\in\rr^{m + m + d_A}. \] 
Since $\bm y\in Y_{E, A}$, we have $\bm 0  = E^\top \bmy = E^\top(\bmu^+ - \bmu^-) =  E^\top\bmu^- -E^\top\bmu^+$; we also have that ${\bmu\rast} \geq \bm 0$ since $A^\top\bmy \leq0$.  That says, 
\[
\left[\begin{array}{c|c|c} E^\top & -E^\top & \bm O \\ A^\top & -A^\top & - I \end{array}\right]\begin{bmatrix}\bmu^-\\ \bmu^+\\ - A^\top \bmy \end{bmatrix} = 
\begin{bmatrix} 
E^\top\bmu^- -E^\top\bmu^+ \\ 
A^\top\bmu^- - A^\top\bmu^+ + A^\top \bmy 
\end{bmatrix} = \bmat{\bm 0 \\ \bm0} = \bm0. 
\] 
Since the system~(\ref{eqn: geo farkas dual}) has no solution, we have that 
\[- \inner{\bmy, \bmb} = \inner{-\bmy, \bmb} = \inner{ \bmu^-, \bmb} - \inner{\bmu^+, \bmb}  =  \inner{
\begin{bmatrix}\bmu^-\\ \bmu^+\\ - A^\top \bmy  \end{bmatrix},
\begin{bmatrix} \bmb\\ -\bmb\\ \bm0 \end{bmatrix}
} \geq 0 \] and, thus $\inner{\bmy, \bmb}\leq 0$. 
\end{proof}

We now consider the geometric intuition behind Theorem~\ref{farkas geo}. Consider the matrices and their columns 
\[E = [\bme_1 \mid \bme_2 \mid \dots \mid \bme_{d_E}]\quad\text{ and }\quad A = [\bma_1 \mid \bma_2 \mid\dots \mid \bma_{d_A}];\]
the theorem states that the neccessary condition, Equation~\ref{eqn: farkas geo} \[\bmb\in \operatorname{span}(\{\bm e_1, \bm e_2, \dots, \bme_{d_E}\})\cup\cone(\{\bm a_1, \bm a_2, \dots, \bm a_{d_A}\})\] happens if and only if, for every vector $\bm y$ that is orthogonal to every $\bm e_i$ and oppose to every $\bm a_i$, $\bm y$ is also oppose to $\bm b$. 

%Finally, we make sure the closure of a finitely generated cone via Farkas' Lemma. Recall that a finite intersection of closed 

